%; whizzy chapter
% -initex iniptex -latex platex -format platex -bibtex jbibtex -fmt fmt
% 以上 whizzytex を使用する場合の設定。

%     Tokyo Debian Meeting resources
%     Copyright (C) 2006 Junichi Uekawa

%     This program is free software; you can redistribute it and/or modify
%     it under the terms of the GNU General Public License as published by
%     the Free Software Foundation; either version 2 of the License, or
%     (at your option) any later version.

%     This program is distributed in the hope that it will be useful,
%     but WITHOUT ANY WARRANTY; without even the implied warranty of
%     MERCHANTABILITY or FITNESS FOR A PARTICULAR PURPOSE.  See the
%     GNU General Public License for more details.

%     You should have received a copy of the GNU General Public License
%     along with this program; if not, write to the Free Software
%     Foundation, Inc., 51 Franklin St, Fifth Floor, Boston, MA  02110-1301 USA


%   Pdf作成手順
% dvipdfmx debianmeetingresume200511.dvi
%  preview (shell-command "xpdf debianmeetingresume200511.pdf&")

%%ここからヘッダ開始。

\documentclass[mingoth,a4paper]{jsarticle}
\usepackage[dvipdfmx]{graphicx}
\usepackage{fancybox}
\usepackage{longtable}
\usepackage{ascmac}	% 囲み (screen,itembox)
\usepackage{fancyvrb}   % 囲み Verbatim のために必要
\usepackage[dvipdfmx]{hyperref}
\usepackage{url}

%http://www.naney.org/diki/dk/hyperref.html
%日本語EUC系環境の時
\AtBeginDvi{\special{pdf:tounicode EUC-UCS2}}
%シフトJIS系環境の時
%\AtBeginDvi{\special{pdf:tounicode 90ms-RKSJ-UCS2}}

%% spacing の設定をする。外枠を減らす。
\setlength\headheight{0mm}
\setlength\topmargin{-20mm}
\setlength\headsep{0mm}
\setlength\topskip{3mm}
\setlength\maxdepth{4pt}
\setlength\columnsep{6mm}
\setlength\textheight{252mm}
\setlength\topmargin{-5mm}
\setlength\textwidth{170mm}
\setlength\oddsidemargin{-5mm}
\setlength\evensidemargin{-5mm}

% commandline環境を定義。画面入出力についてはcommandline環境
% で表記する
\newenvironment{commandline}%
{\VerbatimEnvironment
  \begin{Sbox}\begin{minipage}{15cm}\begin{fontsize}{7.3}{7.3} \begin{BVerbatim}}%
{\end{BVerbatim}\end{fontsize}\end{minipage}\end{Sbox}
  \setlength{\fboxsep}{8pt}\fbox{\TheSbox}}

%%% start of santaku
\makeatletter
\newwrite\tf@jqz
\immediate\openout\tf@jqz\jobname.jqz\relax
\makeatother
\newcounter{santakucounter}
\newcommand{\santaku}[5]{%
\addtocounter{santakucounter}{1}

\addtocontents{jqz}{\arabic{santakucounter}. #5\\}
\nopagebreak 問題\arabic{santakucounter}. 
#1\\
\nopagebreak□ A #2\\
\nopagebreak□ B #3\\
\nopagebreak□ C #4
\pagebreak[1]
\hspace{1cm}
\\

}
%%% end of santaku

\newcommand{\emptyspace}{(\underline{\hspace{1cm}})}



\newcommand{\subsubsubsection}[1]{%
\vspace{1zw}{\bf #1}\\}


% sectionをセンタリングする
\makeatletter
  \renewcommand{\section}{\@startsection{section}{1}{\z@}%
    {\Cvs \@plus.5\Cdp \@minus.2\Cdp}% 前アキ
    {.5\Cvs \@plus.3\Cdp}% 後アキ
    {\normalfont\Large\headfont\raggedright\centering}} % style
\makeatother

% section の代わりの環境
\newcommand{\dancersection}[2]{%
\newpage
東京エリアDebian勉強会 2005
\hrule
\vspace{0.5mm}
\hrule
\hfill{}\includegraphics[width=3cm]{image200502/openlogo-nd.eps}\\
\vspace{-4cm}
\begin{center}
  \section{#1}
\end{center}
\hfill{}#2\hspace{3cm}\space\\
\hrule
\hrule
\vspace{1cm}
}

%% for gotom
\newenvironment{gdescription}%  
{%
   \begin{list}{}% 見出し記号／直後の空白を調節
   {%
      \setlength{\itemindent}{0mm}
      \setlength{\leftmargin}{45mm}%  左のインデント
      \setlength{\rightmargin}{0zw}% 右のインデント
      \setlength{\labelsep}{4mm}%    黒丸と説明文の間
      \setlength{\labelwidth}{4cm}%  ラベルの幅
      \setlength{\itemsep}{0em}%     項目ごとの改行幅
      \setlength{\parsep}{0cm}%      段落での改行幅
      \setlength{\listparindent}{0cm}% 段落での一字下り
      \let\makelabel\gdescriptionlabel
   }
}{%
   \end{list}%
}
\newcommand*\gdescriptionlabel[1]{\hspace\labelsep\normalfont\bfseries #1}
%%

\begin{document}

\begin{titlepage}

% 毎月変更する部分, 本文の末尾も修正することをわすれずに
\title{
 第10回 東京エリア Debian 勉強会\\事前資料}
\date{2005年11月12日}
\author{Debian勉強会会場係 上川 純一\thanks{Debian Project Official Developer}} 
\maketitle
\thispagestyle{empty}

\end{titlepage}

\newpage
\tableofcontents

\dancersection{Introduction To Debian 勉強会}{上川 純一}

今月のDebian勉強会へようこそ。
これからDebianのあやしい世界に入るという方も、すでにどっぷりとつかってい
るという方も、月に一回Debianについて語りませんか？

目的として下記の二つを考えています。

\begin{itemize}
 \item メールではよみとれない、もしくはよみとってられないような情報を情
       報共有する場をつくる
 \item まとまっていないDebianを利用する際の情報をまとめて、ある程度の塊と
       して出してみる
\end{itemize}

また、東京にはLinuxの勉強会はたくさんありますので、Debianに限定した勉強
会にします。Linuxの基本的な利用方法などが知りたい方は、他でがんばってくださ
い。
Debianの勉強会ということで究極的には参加者全員がDebian Packageを
がりがりと作りながらスーパーハッカーになれるような姿を妄想しています。

Debianをこれからどうするという能動的な展開への土台としての空間を提供し、
情報の共有をしたい、というのが目的です。
次回は違うこと言ってるかもしれませんが、御容赦を。

\subsection{講師紹介}

\begin{itemize}
 \item{こばやしさん} 毎回DWNの翻訳をやっていただいています
 \item{えとーさん} dpkg-statoverrideについてかたってもらいます
 \item{上川 純一} 宴会の幹事です。
\end{itemize}

\subsection{事前課題紹介}

今回の事前課題は
「最近dpkgとはこう接しています」

というタイトルで200-800文字程度の文章を書いてください。
というものでした。
その課題に対して下記の内容を提出いただきました。

\subsubsection{吉田@板橋さん}

Debian辞典を購入してやっとdebのわかりやすい作り方を理解したので、
dpkg-buildpackageやdebuildと戯れています。
dpkg -cでなぜ含まれないのか頭を悩ませ、
なんとか形になったら、dpkg -iで叩き込んでいます。

その他、dpkg〜、apt〜も大量にあり、コマンドがいろいろで混乱しやすいため、
メモ的なシェルスクリプトを作って
思い出しやすいrpm系コマンドに似せて使っています。

最近のdpkgとの接し方としては、ようやく打ち解けてきた
というところです。


rpm-ql

\begin{commandline}
 #!/bin/bash
 echo 'dpkg -L (パッケージ名) '
\end{commandline}

rpm-qa--last

\begin{commandline}
 #!/bin/bash
 echo 'ls /usr/share/doc -ltr | tac'
\end{commandline}

 rpm-e 

\begin{commandline}
 #!/bin/bash
 echo '完全削除'
 echo 'apt-get --purge remove (Package name) '
 echo '設定ファイルを残して削除'
 echo 'apt-get remove (Package name) '
\end{commandline}

 rpm-qpi 

\begin{commandline}
 #!/bin/bash
 echo 'Search Package name '
 echo 'auto-apt search (File name) '
 echo 'apt-file search (File name) '
 echo '日本語説明'
 echo 'apt-cache show (パッケージ名)'
\end{commandline}

rpm-qi 

\begin{commandline}
 #!/bin/bash
 echo 'English '
 echo 'dpkg -s (Package name) '
 echo '日本語説明'
 echo 'apt-cache show (パッケージ名)'
\end{commandline}

rpm-qpl 

\begin{commandline}
 #!/bin/bash
 echo 'dpkg -c (パッケージファイル名) '
\end{commandline}

rpm-qa 

\begin{commandline}
 #!/bin/bash
 echo 'dpkg -l'
\end{commandline}

 rpm-qR 

\begin{commandline}
 #!/bin/bash
 echo 'apt-cache depends (package名)'
\end{commandline}

rpm-qf 

\begin{commandline}

 #!/bin/bash
 echo 'dpkg -S (コマンド名) | grep -w `which (コマンド名)`'
 echo 'dlocate (コマンド名やファイル名)'
\end{commandline}


\subsubsection{澤田さん}

dpkgとどう接しているかと言われると、
\begin{itemize}
 
 \item  パッケージを入れた後にどんなファイルが含まれてるか見るために-L
 \item  パッケージを入れてるかを確認するために-l
 \item  ファイルがどのパッケージのものかを確認するために-S
\end{itemize}

ですね。-Lでファイル一覧表示→気になったファイルをlvってのをよくやるの
で、それを統合したアプリなんてあるとうれしいのかもしれない。
\subsubsection{中島　清貴さん}

まったくdpkgとは接していない。こんどちゃんと使ってみようと思う。まだ今のところ
、なんの必要にも迫られていないから使わなくても大丈夫のようだ。しかし自分のマシ
ンに何が入ってるのか確認してないので、これで一括表示とかしてみたら良いのかもし
れない。あと何に使えるのかよく分かってない。dpkg on Solarisというのがあるらし
いので今から調べて使ってみようかと思う。それぐらいだろうか？　とりあえず最近は
これぐらいしか接していないから、あまり使うような接点がないので無理に使うしかな
い。　　　


\subsubsection{上川}

以前はemacsでパッケージ管理ができるとよいのではないか、と勘違いして
apt-elというのを開発していました。
しかし、全然便利じゃないのではないかという致命的なことに気づいたので、あ
きらめています。

最近はdpkgは検索ツールとしてしか使っていません。
dpkg -c や、 dpkg -L, dpkg -S をよく使います。

dpkg -lって凶悪だと思いません？
バージョン番号とパッケージ名がしっかりわかるのが一番重要だと思うのですが、
説明文を表示するために容赦なくバージョン番号を切り捨ててくれます。
これは困りますね。


\subsubsection{えとー}

dpkg は検索ツールとして使うことが多いです。
dpkg -l, dpkg -L, dpkg -S は多用しています。

他に使うのって、自作パッケージを dpkg -i で突っ込むくらいになっています。

dpkg, dpkgていうけど皆さん /usr/sbin/dpkg しかあんま興味ないかもしれませんが、
dpkg-statoverride や dpkg-divert, update-alternatives
なんかはフロントエンド作ってみて、
気にしないでも生きていけるけど知ってると便利なこともあるんだなぁと思うようなものが
あるんだなと思って結構興味持っています。

しかし、新dpkgは話はいっぱいあるのだけどなかなか出てこないのでちょっと寂しいですね。


\subsubsection{小林儀匡}

dpkgは普段は専ら検索ツールとして使っています。
主にdpkg -l, dpkg -L, dpkg -Sを使います。
ただ、dpkg -lはパッケージ名やパッケージバージョンのカラムに入る文字数が少なく、
それを超えると切られてしまうのが厄介です。
おかげで、正確なパッケージ名やバージョン番号が知りたいのに分からなかったり、
パイプでgrepに渡して引っ掛けようとしても引っ掛からないことがよくあります。

他にdpkg-reconfigureも珠に用いますが、
実はこれはパッケージとしてはdebconfに属するのだといまdpkg -Lやdpkg -Sを実行して初めて知りました。

%%% trivia quiz
\dancersection{Debian Weekly News trivia quiz}{上川 純一}

ところで、Debian Weekly News (DWN)は読んでいますか？
Debian 界隈でおきていることについて書いているDebian Weekly News.
毎回読んでいるといろいろと分かって来ますが、一人で読んでいても、解説が少
ないので、
意味がわからないところもあるかも知れません。みんなでDWNを読んでみましょう。

漫然と読むだけではおもしろくないので、DWNの記事から出題した以下の質問にこたえてみてください。
後で内容は解説します。

\subsection{2005年44号}

\url{http://www.debian.org/News/weekly/2005/44/}
にある11月1日版です。

\santaku{Nathanael Nerodeがi386のサポートについて何を宣言したか}
{そろそろi386CPUも世界から駆逐できたので、無くす}
{gccがi386サポートを復活させたので、etchでは真のi386マシンでも動くリリー
スが作れるかも知れない}
{i386向けのバイナリをとうとう生成できなくなってしまったので、今後は対応
はしないことになる。}
{B}

\santaku{Jay Berkenbiltがlibtoolの依存関係からパッケージの依存関係を計算するためのスクリプトを作
成したいという宣言を出した際に、問題になるだろうと指摘されたのは}
{multiarchの場合の.laファイルが複数あるケースをどう扱うか}
{思想的にそういうことはしてはならないことになっているので話題にあげるな}
{自動で依存関係を解決するなんてそんな危険なことをしてはならない}
{A}

\santaku{opensslの新しいバージョンがアップロードされたがその影響は}
{opensslに必要だったセキュリティー対応が取り込まれた}
{opensslがより安定したバージョンになった}
{結果として300以上のパッケージをリビルドする必要がある}
{C}

\santaku{berlinuxでは何が展示されていなかったか}
{nokia 770 で動いているDebian}
{Debianで制御している鉄道模型 }
{Debianを使った炊飯器}
{C}

\santaku{IETFのRFCについてSimon Josefssonは何をしようとしてい
るか}
{IETFからRFCを解放する運動を開始したい}
{RFCを利用しなくても世界がまわるようにあたらしい規格システムを提案、普及
推進}
{RFCのライセンスをフリーソフトウェアに使いやすいように変更するための署名
運動}
{C}

\santaku{openafs-module-sourceの機構で問題になっているのは何か}
{インストールしても動かないため、それをごまかすための機構}
{ユーザに使い方を説明するためのマニュアルが大きすぎること}
{アップグレードする際にモジュールを自動で構築するように求める機構}
{C}

\santaku{自動テストの結果を上流に還元する方法についてDaniel Jacobwitzは
何を提案したか}
{パッケージのビルド中に結果を出力するようにする}
{結果ログを上流に自動でメール送信するように設定する}
{頑張って手動で確認する}
{A}

\santaku{Debian packageで、postinstの処理がパッケージ毎に実行されてしま
うために遅い場合がある。
その処理を複数パッケージ分まとめて実行する方法にはなにがあるか}
{cygwinのsetup.exeを利用する}
{そんな方法はない。}
{zopeはaptのpost-installを利用している。}
{C}

\santaku{gnomeメタパッケージがgnome-gamesに依存していることに出た苦情は}
{政府機関ではゲームが禁止されているのでインストールされると困る。}
{gnome-gamesは大きすぎる}
{gnome-gamesはおもしろくないので、もっとおもしろいゲームを提供すべきだ。}
{A}

\subsection{2005年45号}
\url{http://www.debian.org/News/weekly/2005/45/}
にある11月8日版です。

\santaku{Linux-Info-Tagのブースではネットワーク障害があったが、どうやって対応し
たか}
{ノートパソコンの電源さえあればみんなハッピーだった}
{ノートパソコンにDebianのミラーがはいっていたので、ネットワークは特に問
題ではなかった。}
{通信衛星を利用してネットワークを仮稼働させたので問題にはならなかった}
{B}

\santaku{Robert Milanが発表したgingは何をするものか？}
{CDROMから起動するDebian GNU/kFreeBSD}
{CDROMから起動するWindows}
{CDROMから起動するHurd}
{A}

\santaku{ 次回の debconf の開催期間はいつか}
{2005年5月14日から22日}
{2007年5月14日から22日}
{2006年5月14日から22日}
{C}

\santaku{opensslを利用しているGPLのプログラムをgnutlsを利用するようにするには}
{気合いで書き直す}
{opensslを使い続けるほうが利点があるので、GPLのプログラムのライセンスを
変えてしまう}
{gnutls/openssl.hという互換レイヤーがあるのでそれを利用すればよい}
{C}

\santaku{\url{http://popcon.debian.org/}のトップページで見れない情報は何か}
{最近各アーキテクチャにおいてのpopconの利用者が急激に増えている}
{popconのバージョンがリリースされた際にインストールされている数がどう遷
移するか}
{どのパッケージが一番人気があるか}
{C}


\dancersection{最近のDebian関連のミーティング報告}{上川 純一}

\subsection{東京エリアDebian勉強会9回目報告}

前回開催した第9回目の勉強会の報告をします。

当日のタイムテーブルは下記でした
\begin{itemize}
 \item 18:10- quiz
 \item 18:30- こたえあわせ
 \item 19:00- 休憩
 \item 19:10- たるさん
 \item 20:10- 上川
 \item 21:00- 宴会
\end{itemize}

今回はdebbugsについての熱い話しを展開しました。
今回の参加人数は登録者が11名くらいで、実際に参加したのが9名くらいでした。

DWN quizに関しては、今回は小林さんが1問不正解で最高点数でした。

たるいしさんがapt-listbugsについて説明しました。
昨年のAsia Debian Mini Confで発表した内容を説明して、
実はそのころのTODOは進捗していない、ということを説明していました。
osdn.debian.or.jp でミラーしているのですが、rsync もとはもともと master.debian.orgだったが、
同期できていないというバグ報告があって気づいて、 merkel.debian.org に切替えたという話しがでていました。
資料にあるグラフは２００３年９月から２００４年１０月のもので現在はすでに
一年経過しているため、おそらく７０００IPアドレスからの利用があるのだろう、
と予想していました。
また、中国で発表したときに、使ってみたら、バグの情報を取得するのに何分もかかり、
ミラーサーバ必要だ、という話になった、がいまだになにもできていないとか、
RSSViewerを使えば、自分のマシンに今入っているパッケージのセキュリティー
関連のバグだけを見る、ということができる。
実は便利かもしれない、とか。
現状のコードのメンテナンスにあきてきたのでどうせなら書き直したいなぁという宣言もでていました。

上川が、Anthony TownsがFinlandのdebconfで発表していた内容と、
その後にその発表に触発されて実装されたDebbugsの新機能について話しました。
おそらくDebbugsの仕様について日本語で記述した資料はこれが初めてなのではないでしょうか。
新しい機能がいろいろと実装されており、apt-listbugsでも利用できそうな情報もあるので、
たるいしさんの今後のハックに期待です。

\subsection{東京エリアDebian勉強会9.5回目報告}

大阪で実施した第9.5回目の勉強会の報告をします。

	  10月の関西にての第9.5回Debian勉強会、関西出張版を実施しました。
	  今回はバグレポートの仕方についての話を展開しました。
	  今回の参加人数は40人くらいでした。

	    DWN quizは、今回はえとーさんが司会をしてすすめました。
	    会場の人々は意外にもあまりDWNをよんでいなかったようなので
	    難しい問題が多かったようです。

	    やまねさんがreportbugのしかたについて語りました。
	    まず、バグ報告のなんたるかを解説し、
	    Debianはバグを適正に報告することによって改善されている、ということを説明しました。
	    会場からは、このバグはどうだ、という話題がでました。
	    パッチをつけてバグ報告をするのはパッチをつけないよりはよい。
	    後藤さんがいうには、
	    過去の経験から、バグ報告についているパッチはそれだけでは動かない物が
	    多いので、そのまま適用して終了という場合はすくないのだが、
	    それなりに修正することに対してのモチベーションにつながるし、
	    どういう方向に改変したら修正となるのかというのが明確になるので、
	    好ましい、ということでした。
	    例として、上川が昔バグ報告でファイルしたlvの2GB以上のファイルが扱えないというバグ報告に関して、
	    LARGEFILEオプションでコンパイルしたらよいというパッチをつけていたのですが、
	    後藤さん曰く、LARGEFILEオプションだけではダメで、
	    ftellの代わりにftelloを使う必要があるのです。
	    そのために多くの変更が必要となってしまうため、
	    上川のなげた一行のパッチの適用だけでは問題は解決できません、
	    ということを説明してくれました。

	    その後、上川が、debbugsの構造について解説しました。
	    debbugsのデータファイルの構造や、プログラムの構成についてのおおざっぱな
	    解説をしました。
	    会場にdebbugsの利用方法をきいたところ、
	    ダイレクトにCGIを使っている人は意外と少なく、
	    qa.debian.orgを利用しているメンテナが
	    たくさん居ました。
	    上川の
	      qa.debian.orgのページ
	    \url{http://qa.debian.org/developer.php?login=dancer} 
	    を例にして見てみるとわかりますが、
	    メンテナの管理しているパッケージについての全体像が簡単に把握でき、
	    バグレポートの様子をおおまかに把握することができます。

	    会場の人から、この勉強会はディストリビューションのテストをしている人だけを対象にしているのですか、
	    Debianのインストールの仕方についての話題などはないのですか、
	    たとえば、RAIDを使うのだったらこのパッチが必要だ、という話題などは無いのか
	    ときかれてしまいました。
	    ディストリビューションのテストをしている人達を対象としているというより、
	    現状としては参加している人の多くがDebianの開発に参加している人ですので
	    話題が開発者よりにすこしよっています。
	    インストールの話題については雑誌などを参照にしてやってしまえばよいだろうし、
	    RAIDなどは特にパッチなどは必要無い、という回答をしました。
	    今後要望が多いようであればインストール大会みたいなものも
	    開催することも可能なのですが、
	    いかがでしょうか。

	    最後に武藤さんが司会でPGPキーサインパーティーを開催しました。
	    27人が列をつくって、公開鍵とIDの確認を実施しました。
	    その後、CACertの認証も実施しました。

	    宴会は近所の焼き鳥やさん「とり蔵」で開催。
	    おいしい鳥料理をたくさんいただきました。


\dancersection{Debianとステータス}{えとー}
\label{sec:eto}

\subsection{statoverrideとは？}

ステータス情報を操作するインターフェースファイル、ディレクトリ、
デバイスファイルなどdpkgが扱えるファイルシステムオブジェクトならなんでも扱うことができる。

\subsection{歴史}

dpkg-statoverride の前には dh\_suidregister というのがあった。
その suidregister の幾つかの問題を解決したものとなっている。
2000年のことなのですでに5年以上の歴史のある古いものになっているが今だ有効な仕組み。
\footnote{参考リンク:\url{http://lists.debian.org/debian-dpkg/2000/06/msg00015.html}}

\subsection{どんな風に使われている？}

  \subsubsection{メンテナ編}

    まずは身近なpostfixで用いられているメンテナスクリプトから抜粋してみましょう。

postinst

\begin{commandline}
1 dpkg-statoverride --remove \$POSTDROP >/dev/null 2>&1 || true
2 dpkg-statoverride --remove /var/spool/postfix/public >/dev/null 2>&1 || true
3 dpkg-statoverride --remove /usr/sbin/postqueue >/dev/null 2>&1 || true 
4 dpkg-statoverride --update --add root postdrop 02555 \$POSTDROP
5 dpkg-statoverride --update --add postfix postdrop 02710 /var/spool/postfix/public
6 dpkg-statoverride --update --add root postdrop 02555 /usr/sbin/postqueue
\end{commandline}

postrm  

\begin{commandline}
1 dpkg-statoverride --remove /usr/sbin/postdrop >/dev/null 2>&1 || true
2 dpkg-statoverride --remove /var/spool/postfix/public >/dev/null 2>&1 || true
3 dpkg-statoverride --remove /usr/sbin/postqueue >/dev/null 2>&1 || true
\end{commandline}

  \begin{enumerate}

  \item dpkg-statoverride --remove /usr/sbin/postdrop >/dev/null 2>\&1 || true

        /usr/sbin/postdrop のstatoverrideの設定を削除する

  \item dpkg-statoverride --remove /var/spool/postfix/public >/dev/null 2>\&1 || true

        /var/spool/postfix/public のstatoverrideの設定を削除する

  \item dpkg-statoverride --remove /usr/sbin/postqueue >/dev/null 2>\&1 || true

        /usr/sbin/postqueue のstatoverrideの設定を削除する

  \item dpkg-statoverride --update --add root postdrop 02555 \$POSTDROP

        ユーザ「root」グループ「postdrop」で実行時にグループIDのユーザで起動、読み取りと実行をオーナユーザ、
        オーナグループ、その他に許可したものを/usr/sbin/postdropに付与する。

  \item dpkg-statoverride --update --add postfix postdrop 02710 /var/spool/postfix/public

        ユーザ postfix グループ postdrop で実行時にグループIDを指定した、オーナユーザには
        読み、書き、実行を、オーナグループには実行を、その他にはなにも出来ない権限を
        /var/spool/postfix/pubulicに付与する。

  \item dpkg-statoverride --update --add root postdrop 02555 /usr/sbin/postqueue

        ユーザ root グループ postdrop で実行時にグループIDを指定した、読み取りと実行をオーナユーザ、
        オーナグループ、その他に許可したものを/usr/sbin/postqueue に付与する。

  \end{enumerate}

  以上のように、プログラムにデフォルトのパーミッションを付与したい場合に用いることができる。

  パッケージメンテナがこれを利用したほうがよいのには理由がある。

  debianのポリシーマニュアルでは標準のパーミッションなどが規定されている、しかし、規定されているものでは
  却って不便になってしまうものがあり、ポリシーから離れてパーミッションなどを変更する場合に
  あとからユーザがポリシーにない状態の確認を容易にするために用いるインターフェースとして利便性のためである。

  別の例を上げてみよう。

  \begin{Verbatim}[frame=single]
  for i in /usr/bin/foo /usr/sbin/bar
  do
    if ! dpkg-statoverride --list $i >/dev/null
    then
      dpkg-statoverride --update --add sysuser root 4755 $i
    fi
  done
  \end{Verbatim}

  \begin{itemize}
    \item [解説] 
    \begin{itemize}
      \item /usr/bin/foo や /usr/sbin/bar に statoverride の設定がされていなければ
            ユーザ「sysuser」、グループ「root」、パーミッション「4755」を設定する。
    \end{itemize}
  \end{itemize}

  第二にpostinstなどでchmod,chownなどを使う場合でもユーザの都合でオーナーやパーミッションを変更させることが
  ありうるが、たとえ変更したとしてもパッケージの再インストール時やアップグレード時などにパッケージ標準の
  オーナとパーミッションに上書きされてしまう可能性がある。
  これを防ぐためにパッケージ標準のオーナ及びパーミッションを設定可能にしておくのがよいかと思われる。

  パーミッションをごたごた言うBTSを減らすためにも statoverride は活用しましょう！

  \subsubsection{ユーザ編}

  サーバ管理者などはパッケージにより提供されているオーナやパーミッションでは目的が達成できない場合がある、
  アップデートの少ないstableを使ったとしてもセキュリティアップデートなどでupdateがかかった場合にどうしても
  上書きされてしまい、初期の値に戻されてしまうことになる。
  これを回避するためには、dpkg-statoverrideを使うことによりパッケージの上書きを回避してユーザ設定の
  オーナとパーミッションをアップグレードしても無関係に用いることができるようになる。
  
  GUI環境がある場合は dsys を使ってみてください。

\subsection{使い方}

  \subsection{一般的な用法}

  \begin{enumerate}
  \item ステータスの追加

    \# dpkg-statoverride --add ユーザ名 グループ名 パーミッション ファイル名

  \item ステータスの即時追加]
  
    \# dpkg-statoverride --update --add ユーザ名 グループ名 パーミッション ファイル名

  \item ステータスの削除

    \# dpkg-statoverride --remove ファイル名   

  \item ステータスの変更 

     \# dpkg-statoverride --remove ファイル名 
     \# dpkg-statoverride --update --add ユーザ名 グループ名 パーミッション ファイル名

  \item ステータスの確認
     
     \# dpkg-statoverride --list パターン

  \end{enumerate}


  \subsection{マニアックな用法}

  --admindir オプションを使いホームディレクトリの .（ドット）ファイルのパーミッションを制御する。

  \subsection{パッケージアップグレード}

\begin{enumerate}
  \item postinst などで statoverride が使われているものを管理者などが statoverride の設定を変更した場合

    postinst などで設定を --remove してから --add している場合などにはシステム管理者が statoverride を
    設定しても postinst などの設定に書き換えられてしまう。(postfixなど)

    postinst などで --add のみの場合は管理者などが設定した statoverride の設定になる

  \item postinst などで statoverride をセットされているが管理者が chown、chmod などをした場合

    statoverrideの設定に従ってパーミションが変更される

  \item postinst などで chown や chmod しているものを管理者などが statoverride の設定を変更した場合

    システム管理者が statoverride を設定してもパッケージの postinst などの chown、chmodで書き換えられてしまう。

  \item postinst などでは statoverride も chonw、chmod も使われていないものを statoverride の設定を行なった場合

    statoverride の設定に従ってファイルのパーミッションが変更される

  \item postinst などで chown や chmod しているものを管理者などが chown や chmod で設定を変更した場合

    システム管理者が chown や chmod で設定してもパッケージの postinst などの chown、chmodで書き換えられてしまう。

\end{enumerate}

\subsection{おわりに}

  statoverride は alternatives と比べると単純な機能ではシンプルなものだし知名度は更に低いが、
  メンテナもユーザも楽になれるツールなので是非活用していただきたいと思う。


\dancersection{Debian Weekly News翻訳フロー}{小林儀匡}
\label{sec:kobayashi}

\subsection{はじめに}

Debian Weekly News (DWN) という、
Debianコミュニティのために毎週発行されるニュースレターを御存知でしょうか。
Debian通になるためには読むのが必須とされている (?) ニュースです\footnote{もちろん、
きちんとすべて頭に入れ、忘れないよう毎日読み返す、なんてことはしなくてかまわないでしょう。
ただ、世界に広がるDebianコミュニティにはどんな人がいて、どんな活動をしているか、
Debianがどのような人々の日常的な保守作業や議論を通して作られているかは、
一ユーザとして知っておいたほうがよいように思うので、
目は軽く通しておいたほうがよいと思います。あくまで私見ですが。}。
Debian JPのdebian-usersメーリングリスト\footnote{\url{debian-users@debian.or.jp}。}を購読しているかたなら、
最近では毎週のように翻訳したものが流れてくるので知っているでしょう。
このDWN翻訳作業は、以前は今井伸広さんが一人でされていましたが
2005年6月ごろからDDTP\footnote{Debian Description Translation Project。Debianパッケージ説明文の翻訳をするプロジェクト。}日本語チームコーディネータの田村一平さん、
そして小林がチームを組んでおこなっています。
ここでは、その翻訳作業の流れについて説明します。

\subsection{Debian Weekly Newsとは}

DWN はDebianコミュニティのための週刊ニュースレターで、
UTCで毎週火曜18:45ごろ (JSTで毎週水曜3:45ごろ) に発行されます。
「週刊」というだけあって基本的に毎週発行されますが、発行されないことも年に数回あります。
発行形態はメーリングリストとウェブページの2種類あり、
前者は本家のdebian-newsメーリングリスト\footnote{\url{debian-news@lists.debian.org}。}で購読できます
(もちろんこのメーリングリストにはDWNの他に不定期のニュースも珠に流れてきます)。
また後者は\url{http://www.debian.org/News/weekly/}で参照できます。
様々なメーリングリストでの話題やイベントについて書かれているので、
読んでおくとDebian界隈のニュースに (無駄に) 詳しくなれます。

DWNの編集は最近ではほとんどMartin 'Joey' Schulze一人の作業に委ねられています。
しかしフッタに書かれる編集者情報を見る限り、
2号に1号くらいは他の人が助けているようです。

ちなみに創刊号は1999年1月4日に出され、その編集者はJoey Hessでした。
記念すべき第1号のeditorialによれば、
Linux Weekly News\footnote{\url{http://lwn.net/}。}を真似て作られたようです。
またeditorialの次に書かれた記事第1号のタイトルは `RMS is using Debian.' です。

\subsection{ファイル形式 -- WML}

ウェブページ版のDWNはDebianのサイトの一部なので、
他のページと同じCVSリポジトリ\footnote{リポジトリURLは\url{:pserver:anonymous@cvs.debian.org:/cvs/webwml}で、
ウェブインタフェースも\url{http://cvs.debian.org/?root=webwml}で利用できます。
この中の\texttt{webwml}ディレクトリに\url{http://www.debian.org/}用のwmlファイルやそれからウェブページをビルドするツールなどが含まれています。
例えばDWN 2005年第44号英語版は\texttt{webwml/english/News/weekly/2005/44/index.wml}にあり、
その日本語版は\texttt{webwml/japanese/News/weekly/2005/44/index.wml}にあります。}、
そして同じWMLというファイル形式を利用しています。

ファイル形式には他のページと同様にWMLが用いられています。
WMLとはウェブサイトメタ言語 (web site meta language) のことで、
Debianでは\verb|wml|パッケージとして提供されています。
簡単に言ってしまえばHTMLに命令を混ぜたような言語で、
例えばウェブページのコンテンツに定型的なヘッダ・フッタ、
さらにHTMLのheadを加えたりするのに用いられています。
また、条件分岐を用いた、
「○○のリリース前にはこのコンテンツを表示し、リリース後にはこのコンテンツを表示する」
というような表示の切り替えや、
Debianのウェブページの翻訳版がオリジナル版のどのバージョンに基いているかを管理するのにも用いられています。
WMLについてより詳しく知りたければ\url{http://www.debian.org/devel/website/using_wml}を参照してください\footnote{ちなみに筆者は、WMLがDebian以外にどこで使われているか知りません}。

\subsubsection{DWNのWML}

Debian Weekly News の場合だとWMLの形式は決まっています。
以下では (わざわざ説明する必要もない、見れば分かるものだと思いますが) DWN 2005年第44号のものを例にとって説明します。

まず最初に命令が決ます。
ここには、発行日時・そしてアーカイブのページに表示されるサマリ (主だった記事のキーワードのリスト) が書かれます。
その後にCVSのIdが書かれています。
\begin{verbatim}
#use wml::debian::weeklynews::header PUBDATE="2005-11-01" SUMMARY="Dependencies, OpenSSL, Berlinux, RFCs, Kernel, Packaging, GTK, GNOME"
# $Id$
\end{verbatim}

その後で、
`Welcome to this year's \verb|nth| issue of DWN, the weekly newsletter for the Debian community.'
で始まるeditorialが来ます。

\begin{verbatim}
<p>Welcome to this year's 44th issue of DWN, the weekly newsletter for the
Debian community.  Nathanael Nerode <a
href="http://lists.debian.org/debian-devel/2005/10/msg00388.html">reported</a>
that current GCC versions support the old i386 processor again and hence
Debian could retain i386 compatibility in the upcoming <a
href="$(HOME)/releases/etch/">etch release</a>.</p>
\end{verbatim}

あとは普通の記事の繰り返しです。

\begin{verbatim}
<p><strong>Calculating Development Package Dependencies.</strong> Jay
Berkenbilt <a
href="http://lists.debian.org/debian-devel/2005/10/msg00184.html">\
proposed</a> to work on a <a
href="http://packages.debian.org/debhelper">debhelper</a> script that helps
calculating <a href="http://packages.debian.org/libtool">libtool</a>
dependencies for development packages.  Goswin von Brederlow <a
href="http://lists.debian.org/debian-devel/2005/10/msg00519.html">pointed
out</a> that with <a href="http://raw.no/debian/amd64-multiarch-2">\
multiarch</a> there may be concurrent <code>.la</code> files to handle.  No
consensus in favour of such a script was reached.  Junichi Uekawa (&#19978;&#24029; &#32020;&#19968;)
<a
href="http://lists.debian.org/debian-devel/2005/10/msg00316.html">\
mentioned</a> the <a href="http://packages.debian.org/d-shlibs">d-shlibs</a>
package that contains scripts to support the maintainer in this regard.</p>
\end{verbatim}

以上がDWNの記事部分です。

記事部分のあとには、箇条書で書けるような様々な情報のコーナーが続きます。
主に前号の発行以降に変化があったパッケージについての情報で、
最近載っているのは
\begin{description}
 \item[`Security Updates.' (「セキュリティ上の更新。」)]
   DSA番号・パッケージ名・脆弱性についての記述) のリスト。
 \item[`New or Noteworthy Packages.' (「新規または注目すべきパッケージ。」)]
   パッケージ名・パッケージ説明文のリスト。
 \item[`Orphaned Packages.' (「みなしご化されたパッケージ。」)]
   パッケージ名・パッケージ説明文・debbugsのバグ番号のリスト。
 \item[`Removed Packages.' (「削除されたパッケージ。」)]
   パッケージ名・パッケージ説明文・パッケージ削除理由のリスト。
\end{description}
の4コーナーですが、このうち「削除されたパッケージ。」は比較的新しい項目です。
去年は、
Debian Package a Day's Journal\footnote{\url{http://www.livejournal.com/users/debaday/}。}で紹介されたパッケージのリストを含む
`Debian Packages introduced last Week. Ever' (「先週紹介された Debian パッケージ。」) というコーナーがあったのですが、
これは、ウェブページが更新されなくなった (さすがにネタが尽きた?) ために廃止されたようです。
以下にコーナーの例を示します。

\begin{verbatim}
<p><strong>Security Updates.</strong> You know the drill.  Please make sure
that you update your systems if you have any of these packages installed.</p>

<ul>
<li>DSA 872: <a href="$(HOME)/security/2005/dsa-872">koffice</a> --
    Arbitrary code execution.
<li>DSA 873: <a href="$(HOME)/security/2005/dsa-873">net-snmp</a> --
    Denial of service.
<li>DSA 874: <a href="$(HOME)/security/2005/dsa-874">lynx</a> --
    Arbitrary code execution.
<li>DSA 875: <a href="$(HOME)/security/2005/dsa-875">openssl094</a> --
    Cryptographic weakness.
<li>DSA 876: <a href="$(HOME)/security/2005/dsa-876">lynx-ssl</a> --
    Arbitrary code execution.
<li>DSA 877: <a href="$(HOME)/security/2005/dsa-877">gnump3d</a> --
    Several vulnerabilities.
<li>DSA 878: <a href="$(HOME)/security/2005/dsa-878">netpbm-free</a> --
    Arbitrary code execution.
</ul>
\end{verbatim}

最後に定型メッセージが入ります。

\begin{verbatim}
<p><strong>Want to continue reading DWN?</strong> Please help us create this
newsletter.  We still need more volunteer writers who watch the Debian
community and report about what is going on.  Please see the <a
href="$(HOME)/News/weekly/contributing">contributing page</a> to find out how
to help.  We're looking forward to receiving your mail at <a
href="mailto:dwn@debian.org">dwn@debian.org</a>.</p>
\end{verbatim}

そしてフッタ用のWML命令が入ります。
編集者の名前を書くのに使われます。

\begin{verbatim}
#use wml::debian::weeklynews::footer editor="Martin 'Joey' Schulze"
\end{verbatim}

\subsection{翻訳作業}

翻訳作業はだいたい以下のような流れになっています。

\begin{description}
 \item[週末] Martin `Joey' Schulzeがプライベートな草稿\footnote{\url{http://www.infodrom.org/~joey/Writing/DWN/}。}に記事を追加する。
 \item[水曜未明] オリジナル版がリリースされる。
 \item[水曜] 今井さんが翻訳チームの田村さん・小林に作業を割り振る。
 \item[水曜〜日曜] 翻訳チームメンバーの3人が個々に翻訳し、
   Debian JPのdebian-wwwメーリングリスト\footnote{\url{debian-www@debian.or.jp}。}にそれを投稿して査読をお願いする。
   それに対して杉山さん・武井さんなどが査読をしてくださる。
 \item[月曜〜火曜] (基本的に今井さんが) debian-usersに日本語版をリリースする。
\end{description}

\subsubsection{Martin `Joey' Schulzeのオリジナル版リリース作業}

Martin `Joey' Schulzeは、
オリジナル版をリリースするために独自のリポジトリ
\footnote{\url{:pserver:anonymous@cvs.infodrom.org:/var/cvs/infodrom.org}の
\texttt{public\underline{ }html/src/Writing/DWN}。}を用いています。
彼はそれに記事を追加していき、
UTCで月曜日の夕方に、「ごく親しい関係者」向けにDWNのプレビューをリリースします。
彼はこのプレビューを\verb|dwn@debian.org|に送り、DWNチームのレビューを受けます。
またプレビューは、
早目に翻訳にとりかかりたい人のために\verb|dwn-trans@debian.org|にも送られます。
こうして、
いくつかの言語ではオリジナル版の正式なリリースと同時に翻訳版をリリースすることができるようになっています。
日本語版は現在、正式リリースを待って作業を開始しています。

\subsubsection{割り振り}

翻訳チームメンバー3人への作業の割り振りは基本的に今井さんが行います。
チーム制を始めたばかりのころは「記事部分前半」・「記事部分後半」・「それ以外 (様々なコーナー)」という分け方をしていました。
しかし、記事というのは流れがあるので訳していてそれなりに楽しいのですが、
「それ以外 (様々なコーナー)」は、
脆弱性・パッケージ削除理由といった定型文句が多数を占めるものがある一方で、
流れがなくて訳しづらいパッケージ説明文もあります。
(どちらにしろローテーション制なのですが) それだとやりがいに差が出るためか\footnote{本当の理由は知らないので、もしかしたら違うかもしれません。}、
最近では「記事部分の1/3とどれかのコーナー」という分け方になっているようです。

\subsubsection{翻訳・査読依頼}

翻訳作業は、
CVSリポジトリからオリジナル版のDWNのwmlファイルをコピーして始めます。
このときファイルのコピーには、
同じリポジトリに入っている\verb|webwml/copypage.pl|を用います。
というのも、このPerlスクリプトはコピー元のオリジナル版wmlファイルに含まれるCVSのIdキーワードの値から翻訳バージョンチェック用のWML命令を作り出してくれるからです。
またこのスクリプトは、
オリジナル版wmlファイルに含まれているLatin-1文字コードの文字を適切な文字実体参照に置き換え、
ASCIIだけを含むファイルにしてくれます。
ウェブインタフェースを使ったダウンロードや、
\verb|cp|コマンドによる生のコピーでは、それらがおこなわれません。

実際の翻訳作業では、査読のしやすさを考えて、
オリジナルの記事・リスト項目を残し、その後ろに翻訳をつけていきます。
「この文の意味がよく分からない」「この単語はこう解釈した」といったメモを査読者に残したい場合は、
「\verb|#|」でコメントアウトしたものをオリジナルと翻訳の間に挟んで残します。

\subsubsection{日本語版リリース作業}

リリース作業は次のような手順でおこないます。

\begin{enumerate}
 \item Debian JP debian-wwwメーリングリストに送られた3人の翻訳をwmlファイルにマージする。
 \item 査読をしてくださった方々のコメントを反映させる修正をおこなう。
 \item w3mで表示させ、査読でも指摘されずに残った明らかなミスがないか、
   自分の目で簡単にチェックする。
 \item w3mの表示を見ながら、
   全角文字同士や全角文字とASCIIの間に不必要な空白ができていないか、
   またはASCIIの周りに空白がきちんとあるかを調べ、
   必要に応じて改行位置やタグの位置を変える。
   詳しくは\ref{subsec:translation-care}を参照のこと。
 \item 最下行にあるフッタ用WML命令に、
   \verb|editor|フィールドとは別に\verb|translator|フィールドをつけることができるので、
   そこに翻訳者の名前を連ねる。
 \item 査読しやすいように残してあったオリジナルの記事・リスト項目と翻訳時のメモを消す (これでウェブ用翻訳版wmlファイルは完成となる)。
 \item 今井さんのRubyスクリプトで、
   wmlファイルをDebian JP debian-usersメーリングリスト投稿用プレインテキストに変換する。
 \item \verb|&ouml;|などの文字実体参照をw3mなどの出力に合わせて変換する。
 \item Emacsの\verb|auto-fill-mode|を利用して適当なところで改行する。
   ただし、
   リンクの代わりとなるURLリスト内の項目番号を示す「\verb|[1]|」のような文字列の前で改行されてしまった場合は、
   それを手動で前の行末に戻す。
 \item 以前のdebian-usersへの投稿を元に体裁を整える。
 \item コメントを加える。
   大変な査読作業をしてくださった方々への謝辞を忘れないようにする。
 \item 完成した投稿用プレインテキストをdebian-usersに送信する。
 \item 最後にウェブ用翻訳版wmlファイルをdebian-wwwに送信し、
   コミットをお願いする (もちろん、コミット権のある今井さんが作業者の場合は直接コミットとなる)。
\end{enumerate}

\subsection{翻訳作業時の注意点}
\label{subsec:translation-care}

翻訳作業時には以下のような点に注意を払います。

\begin{itemize}
 \item ウェブページの見栄えを意識して、ASCIIと全角文字の間には空白を入れる。
 \item 変なところで改行すると、
   テキストブラウザで見たときに全角文字どうしの間に空白が入ってしまうことがある。
  それを防ぐため改行位置には注意を払い、たとえば次のような方法をとる。
 \begin{itemize}
  \item 見かけに反映されることのない、タグの内側での改行を多用する。
  \item 特にタグ直後の改行には注意する。
  \item 結果的にスペースを伴うような改行がどうしても必要なときは、
    行末に\verb|\|を入れる。
 \end{itemize}
 \item 改行位置に自信がないときはw3mで見てみる\footnote{\texttt{-T text/html}とオプションを指定すれば、wmlファイルをHTMLファイルとして見ることができます。}。
\end{itemize}

\subsection{今後の課題}

今後の課題としては以下のようなものがあるでしょう。
中にはDWNの翻訳に留まらず他の翻訳にも関係するものもあります。

\begin{description}
 \item[定型文句の翻訳の自動化]
   定型文句を翻訳するときには、訳語や表現形式を合わせるために、
   その文句が使われている過去の翻訳を探してcopy and pasteするという作業が必要になります。
   この作業は、
   オリジナル版と翻訳版を行き来して該当部分を探し出さなければならないので、
   案外面倒なものです。
   対訳をデータベース化し、
   そのデータを用いて翻訳できる部分は作業前に予め機械的に翻訳できるようにしておくと、
   もう少し効率よく作業を進められる気がします。
 \item[対訳表の保守]
   Debian JPに対訳表\footnote{HTML版が\url{http://www.debian.or.jp/Documents/trans_table/trans_table.html}にあり、
	    dict形式のものも
	    \url{http://www.debian.or.jp/devel/doc/about-trans-table.html}
	    で手に入ります。} があるのですが、現在は更新されていません。
   APT関連、Debianプロジェクト関連、
   あるいはその他のDebianでよく使われる用語の対訳表\footnote{翻訳者を煩わせるサードパーティのアプリケーション名の表記に関するポリシーも一緒に含まれていたらよいと思います。}を整備できたらと思います。
   新規に翻訳作業に携わりたい人にとって、
   対訳が簡単に調べられることはかなり重要でしょう。
\end{description}

\subsection{おまけ: 翻訳作業に有用なツール}

小林が翻訳作業によく用いているのは、
英辞郎\footnote{\url{http://www.eijiro.jp/}。}、
NDTPD\footnote{\url{http://www.sra.co.jp/people/m-kasahr/ndtpd/}。
Debianパッケージは\url{http://packages.debian.org/ndtpd}。}、
Lookup\footnote{\url{http://openlab.ring.gr.jp/lookup/}。
Debianパッケージは\url{http://packages.debian.org/lookup-el}。}の組み合わせです。
英辞郎のデータをFreePWING\footnote{\url{http://www.sra.co.jp/people/m-kasahr/freepwing/}。
Debianパッケージは\url{http://packages.debian.org/freepwing}。}で変換して得られたデータに、
Emacs上のクライアントLookupからサーバNDTPDを経由してアクセスします。

また、文書の変更を管理するのに、
バージョン管理システムとしてSubversion\footnote{\url{http://subversion.tigris.org/}。
Debianパッケージは\url{http://packages.debian.org/subversion}。}を利用しています。
差分をunified diffで表示しても分かりにくいときは、
DocDiff\footnote{\url{http://www.kt.rim.or.jp/~hisashim/docdiff/}。
Debianパッケージは\url{http://packages.debian.org/docdiff}。}で表示させると分かりやすくなることがあります。


\newpage 

\dancersection{次回}{}

東京での次回は12月10日土曜日の夜を予定しています。
内容は本日決定予定です。

参加者募集はまた後程。

\newpage

\vspace*{15cm}
\hrule
\vspace{2mm}
\includegraphics[width=2cm]{image200502/openlogo-nd.eps}
\noindent \Large \bf Debian 勉強会資料\\ \\
\noindent \normalfont 2005年11月12日 \hspace{5mm}  初版第1刷発行\\
\noindent \normalfont 東京エリア Debian 勉強会 （編集・印刷・発行）\\
\hrule

\end{document}
